\section{Software Architecture and Implementation Details}
\label{app:software_architecture}

The project is organized into two main directories:
\texttt{src}, which contains the shared numerical
components, and \texttt{Studies}, which contains the
scripts for snapshot generation, POD construction,
Galerkin projection, refinement, and evaluation.

\subsection{Shared numerical modules}

The main source modules are responsible for the
computational grid, initialization, boundary conditions,
diagnostics, and full-order solution.

\paragraph{\texttt{src/grid.py} and
\texttt{src/initialization.py}.}

The function \texttt{create\_grid} generates the
structured grid coordinates. The initialization
module provides the initial temperature field through
\texttt{initialize\_temperature}, which initializes
$T(y)=y$, and constructs the prescribed velocity
components using \texttt{initialize\_velocity\_fields}.

\paragraph{\texttt{src/boundary\_conditions.py}.}

The function \texttt{apply\_temperature\_boundary\_conditions}
implements the four temperature boundary conditions
using ghost cells. It applies the prescribed inlet
and wall temperatures and the zero-gradient outlet
condition.

\paragraph{\texttt{src/diagnostics.py}.}

The function \texttt{build\_rhs} evaluates the
non-temporal terms of the discretized energy equation.
The function \texttt{steady\_state\_residual\_rms}
calculates the RMS residual over the computational
domain. This residual is used to determine convergence
of the full-order solver.

\paragraph{\texttt{src/solvers/direct.py}.}

The full-order solver is implemented through four
main functions:
\begin{itemize}
    \item \texttt{build\_spatial\_matrix}: Constructs
    the sparse five-diagonal spatial matrix. It maps
    physical two-dimensional indices to the
    one-dimensional system and incorporates
    ghost-cell contributions into the matrix
    coefficients.

    \item \texttt{build\_rhs}: Constructs the
    right-hand-side contribution for an individual
    row.

    \item \texttt{build\_forcing\_field}: Applies
    \texttt{build\_rhs} over the computational domain
    to assemble the forcing vector.

    \item \texttt{solve\_with\_direct}: Assembles
    the spatial matrix and forcing vector, factorizes
    $(\mathbf{I}-\Delta t\mathbf{A})$, advances the
    temperature field, and checks convergence using
    the steady-state RMS residual.
\end{itemize}

\subsection{POD construction and Galerkin implementation}

\paragraph{Original snapshot generation and POD basis.}

The script
\texttt{F05\_refinement/f00\_generate\_large\_training\_cases.py}
generates the 100-case training parameter grid. Its
function \texttt{build\_parameter\_grid} creates the
parameter combinations. Each case is then solved
using \texttt{run\_fom\_case}. Converged temperature
fields are converted into snapshot columns, and
their associated simulation information is stored
in \texttt{snapshot\_case\_records}.

The script
\texttt{F05\_refinement/f01\_build\_pod\_basis.py}
loads the snapshot matrix and its case records,
performs the SVD, calculates the POD energy
distribution, and saves the resulting basis.

\paragraph{\texttt{F04\_galerkin\_rom/f01\_operators.py}.}

This module constructs the full-order operator
components required for Galerkin projection.

The function \texttt{\_assemble\_spatial\_matrix}
builds the spatial matrix used to extract the
convection and diffusion operators.

The function \texttt{build\_linear\_operators}
constructs $\mathbf{C}_h$ and $\mathbf{L}_h$.
The function \texttt{build\_viscous\_dissipation\_source}
calculates the parameter-independent source vector
$\boldsymbol{\Phi}_h$ from the prescribed velocity field.

The function \texttt{build\_lifting\_residual}
calculates the discrete lifting contribution, while
\texttt{\_assemble\_forcing\_field} obtains the
affine forcing through the existing full-order
forcing routine.

The function \texttt{verify\_operator\_decomposition}
provides checks for the consistency of the operator
decomposition and lifting-residual formulation.
Finally, \texttt{build\_operator\_components}
collects the resulting matrices and vectors.

\paragraph{\texttt{F04\_galerkin\_rom/f02\_reduced\_operators.py}.}

This module projects the full-order operators onto
the retained POD basis. It constructs the reduced
convection and diffusion matrices and the projected
viscous-dissipation and lifting vectors:
\[
\begin{aligned}
\mathbf{C}_r&=\mathbf{V}_r^T\mathbf{C}_h\mathbf{V}_r,\\
\mathbf{L}_r&=\mathbf{V}_r^T\mathbf{L}_h\mathbf{V}_r,\\
\boldsymbol{\Phi}_r&=\mathbf{V}_r^T\boldsymbol{\Phi}_h,\\
\mathbf{g}_r&=\mathbf{V}_r^T\mathbf{g}_{\mathrm{lift}}.
\end{aligned}
\]

\paragraph{\texttt{F04\_galerkin\_rom/f03\_solve\_reduced.py}.}

This module handles reduced-system assembly,
solution, and temperature reconstruction.

The function \texttt{construct\_reduced\_system}
assembles the reduced matrix and forcing vector.
The function \texttt{solve\_reduced\_system}
calculates the modal coefficients and returns
the reduced-system residual. The function
\texttt{reconstruct\_temperature} recovers the
temperature field from the modal coefficients
and reference temperature.

The main function, \texttt{run\_reduced\_model},
coordinates these operations and saves the
reduced system, modal coefficients, and
reconstructed temperature field.

\paragraph{Case-level ROM evaluation.}

The script
\texttt{F04\_galerkin\_rom/f04\_rom\_for\_case.py}
provides case-level full-order and reduced-order
evaluations. Its main function,
\texttt{run\_rom\_for\_case}, processes the
specified cases, computes their error metrics,
saves the reconstructed fields and records,
and identifies the minimum- and maximum-error
cases for visualization.

The functions \texttt{run\_fom\_case} and
\texttt{run\_rom\_case} provide the corresponding
full-order and reduced-order calculations.
The script
\texttt{F04\_galerkin\_rom/f05\_test\_rom.py}
generates parameter cases and passes them to
the case-level evaluation routine.

\subsection{Snapshot enrichment and independent evaluation}

The refinement workflow is implemented through
the scripts in \texttt{F05\_refinement}.

\paragraph{Mode-count study.}

The script
\texttt{f03\_mode\_count\_study.py}
evaluates the original POD basis over the LHS
screening dataset. It constructs the full-order
operators, evaluates the ROM for each selected
mode count, records the error metrics, and
generates the error-statistics plots and
triangulated error maps.

\paragraph{Hotspot selection and additional snapshots.}

The script
\texttt{f04\_generate\_augmented\_training\_set.py}
identifies high-error candidates and selects
spatially diverse hotspot centres.

The function \texttt{build\_hotspot\_candidates}
ranks the high-error cases separately for each
selected mode count, groups repeated parameter
combinations, and calculates their mean rank
scores, support fractions, and base scores.

The function \texttt{select\_hotspots} applies
the spatial-diversity criterion. The function
\texttt{build\_augmented\_parameter\_grid}
generates the local parameter combinations
around the selected centres and removes
invalid, duplicate, and previously sampled
points.

Finally, \texttt{run\_augmented\_training\_set\_generation}
solves the additional full-order cases and
saves their temperature fields and records.

\paragraph{Augmented POD basis.}

The script
\texttt{f05\_build\_augmented\_pod\_basis.py}
combines the original and additional snapshot
matrices through
\texttt{combine\_snapshot\_datasets}. The
combined matrix is then decomposed using SVD
to obtain and save the augmented POD basis.

\paragraph{Independent test generation.}

The script
\texttt{f06\_generate\_independent\_test\_set.py}
generates the final LHS test dataset. The
function \texttt{load\_training\_parameter\_keys}
loads the parameter combinations used in the
original and augmented training datasets.
The function \texttt{generate\_test\_parameter\_grid}
generates unique test combinations and excludes
previously used training points.

The main function,
\texttt{run\_independent\_test\_set\_generation},
runs the full-order solver for the accepted
test cases and saves their temperature fields
and simulation records.

\paragraph{Original-versus-augmented comparison.}

The script
\texttt{f07\_compare\_original\_and\_augmented\_bases.py}
performs the final paired comparison.

The function \texttt{run\_basis\_comparison}
loads the independent test cases and both POD
bases, evaluates them at the selected mode
counts, and generates the comparison results.

The function \texttt{run\_rom\_evaluations}
applies each basis to exactly the same test
cases. The function \texttt{build\_comparison\_summary}
calculates the aggregate error statistics,
while \texttt{build\_paired\_comparison}
matches the original and augmented results
case by case. The function
\texttt{build\_improvement\_summary}
then calculates the error changes and
improvement statistics.

Finally, the plotting functions generate the
aggregate error comparisons, error-reduction
plots, and parameter-space error maps.

\subsection{Overall execution sequence}

The main computational workflow can be summarized
as follows:
\begin{enumerate}
    \item Generate full-order training cases and
    construct the original POD basis.
    \item Build the discrete Galerkin operators
    and evaluate the original ROM.
    \item Generate the LHS screening dataset and
    perform the mode-count study.
    \item Select hotspots and generate additional
    full-order snapshots.
    \item Construct the augmented POD basis.
    \item Generate the independent test dataset.
    \item Evaluate both POD bases on the same test
    cases and calculate the paired error statistics.
\end{enumerate}

This organization separates the shared numerical
components from the individual study stages and
allows the full-order solver, POD routines, and
Galerkin implementation to be reused throughout
the refinement and evaluation workflow.
